% fig1_data_splits.tex: generated by gmpe-bench; graphics path relative to results/paper/
\begin{figure}[htbp]
\centering
\adjustimage{max width=\textwidth}{figures/fig1_data_splits.pdf}
\caption{Dataset and benchmark splits (full profile: 20,926 records, 575 events, 4,026 stations). (a) Magnitude versus rupture distance by hypocentre region (record counts in the legend). (b) Station $V_{S30}$; dotted lines mark the NEHRP site-class boundaries E|D at 180 m/s, D|C at 360 m/s, C|B at 760 m/s and B|A at 1500 m/s (only boundaries within the data range are drawn). (c) Usable records per period after masking $T > 1/\mathrm{LUF}$. (d) Share of test records whose event or station also occurs in training, averaged over the folds of each split scheme. Split schemes: Rand = record-level random split (leaky reference); Evt = event-grouped 5-fold; Stn = station-grouped 5-fold; Reg = leave-one-region-out; Temp = temporal (train before 2008, test 2008 onward); $M_w{\geq}7$ = magnitude hold-out (test events with $M_w \geq 7$); CL1$\rightarrow$2 = mainshocks (CL1) for training, aftershocks (CL2) for testing.}
\label{fig:fig1_data_splits}
\end{figure}
